\begin{textsample}{POS Dim 1 – vem\_ed\_17 – Score 7.00 – vem\_ed\_17/t078.txt}  \label{ex:f1_pos_134}
A VISÃO DOS DOCENTES: PESQUISA 2º SEMESTRE 2022 \textbf{Todos} os 44 professores \textbf{que} responderam à pesquisa de percepção autorizaram divulgar as respostas para fins de publicações gerenciais e acadêmicas e \textbf{todos} recomendaram um PCI para colegas, como ocorreu nas pesquisas \textbf{anteriores}.
Quase dois terços (64 \%) \textbf{são} veteranos, realizando o PCI pela \textbf{terceira} vez ou mais: isso mostra o amadurecimento das colaborações internacionais nas Fatecs.
Na edição \textbf{anterior} da pesquisa (primeiro semestre de 2022), 49 \% dos professores respondentes realizavam o PCI pela primeira vez.
Aquele \textbf{foi} um \textbf{momento} de captação, por \textbf{parte} da equipe dos PCIs / Cesu, de novos docentes interessados, \textbf{que} mantiveram suas colaborações e aos poucos ganharam experiência internacional.
Essa consolidação se mantém para o \textbf{futuro} próximo: dentre os respondentes do segundo semestre de 2022, 66 \% professores pretendem continuar com mesmo projeto e parceiro, 23 \% continuam com o mesmo parceiro, apenas mudando o projeto, 7 \% vão buscar novo parceiro e apenas 4 \% não desejam realizar outro PCI.
Dos 44 pesquisados, 77,3 \% \textbf{reconhecem} \textbf{que} a \textbf{competência} em língua estrangeira melhorou. \textbf{É} quase unânime (97,8 \%) \textbf{que} os PCIs possibilitam desenvolver novas atividades acadêmicas a \textbf{partir} da experiência do parceiro internacional.
Além disso, 93 \% concordam \textbf{que} os PCIs abrem portas para \textbf{publicar} pesquisas em sua área de atuação e 86,3 \% afirmam \textbf{que} podem elaborar artigos acadêmicos com o parceiro.
Esses dados mantêm a tendência positiva dos anos \textbf{anteriores}.

% matched lemmas: anterior, competência, futuro, momento, parte, partir, publicar, que, reconhecer, ser, terceiro, todo
\end{textsample}
